\section*{Chapitre \ref{ch:b3:milky-way-galaxies} --- La Voie lactée et les galaxies}

\begin{solution}{exo:b3:milky-way-galaxies:1}
(a)~$2\pi R_0 = 2\pi\times\qty{2.523e17}{km} = \qty{1.585e18}{km}$~; à
\qty{229}{km.s^{-1}} cela prend \qty{6.92e15}{s}, soit
\qty{219}{Ma}. (b)~$4540/219 = 20.7$~: environ vingt tours.
\end{solution}

\begin{solution}{exo:b3:milky-way-galaxies:2}
(a)~$\qty{1}{au}/\qty{1}{an} = \num{1.496e8}/\num{3.156e7} =
\qty{4.74}{km.s^{-1}}$. (b)~$v = 4.74\times6.411\times8.178 =
\qty{248.5}{km.s^{-1}}$, compatible avec la valeur
\qty[separate-uncertainty=true]{248 \pm 3}{km.s^{-1}} de la revue. (c)~$248.5 - 11 = \qty{237.5}{km.s^{-1}}$,
environ \qty{8}{km.s^{-1}} au-dessus des \qty{229}{km.s^{-1}} des géantes rouges~:
les deux reposent sur des valeurs différentes de $R_0$ (8.178 et \qty{8.122}{kpc}),
et le \omterm{def:b3:milky-way-galaxies:lsr}{référentiel local au repos} peut lui-même dériver d'environ \qty{15}{km.s^{-1}} par rapport
à une orbite réellement circulaire.
\end{solution}

\begin{solution}{exo:b3:milky-way-galaxies:3}
(a)~$n = 10(1 - 0.6) = 4$~: E4. (b)~$b/a = 1 - 7/10 = 0.3$. (c)~La classe
utilise le rapport d'axes apparent~: un ellipsoïde aplati vu selon son petit
axe présente un contour circulaire et est classé E0.
\end{solution}

\begin{solution}{exo:b3:milky-way-galaxies:4}
$d = 10^{\mu/5 + 1}$~pc~: le LMC \qty{49.6}{kpc}, Andromède
\qty{783}{kpc}, NGC~4258 \qty{7.58}{Mpc}. Une étoile du Nuage a une
parallaxe de $1/\qty{49.6}{kpc} = \qty{0.020}{mas}$, soit la taille de
l'incertitude du meilleur catalogue pour les étoiles brillantes~: une parallaxe isolée ne donne
aucune distance, et le Nuage est situé par d'autres moyens (binaires
à éclipses), à partir desquels s'établit l'échelle des \omterm{def:b3:milky-way-galaxies:cepheid}{céphéides}.
\end{solution}

\begin{solution}{exo:b3:milky-way-galaxies:5}
(a)~$\omega_0 = A - B = \qty{27.2}{km.s^{-1}.kpc^{-1}}$ et $V_0 =
27.2\times8.178 = \qty{222}{km.s^{-1}}$. (b)~$V'_0 = -(A + B) =
\qty{-3.4}{km.s^{-1}.kpc^{-1}}$~: une courbe qui décroît doucement. (c)~Les deux valeurs sont
inférieures à celles des géantes rouges (229 et \qty{-1.7}{km.s^{-1}} par
kiloparsec)~; les auteurs de cette courbe ont trouvé la même chose, avec un $V_0$ environ deux
écarts types plus bas en tenant compte de $R_0$ et de ces \omterm{def:b3:milky-way-galaxies:rotation-curve}{constantes d'Oort}. Les
formules d'Oort supposent des orbites circulaires dans un disque axisymétrique~; des valeurs non nulles de $C$ et
$K$ montrent que le champ de vitesses local est perturbé par la barre et les
bras spiraux, et les $A$ et $B$ locaux portent ces perturbations.
\end{solution}

\begin{solution}{exo:b3:milky-way-galaxies:6}
(a)~$\rho = V^2/(4\pi GR^2) = 229^2/(4\pi\times\num{4.301e-6}\times
8.178^2) = \num{1.45e7}$ masses solaires par kiloparsec cube, soit
\num{0.0145} par parsec cube. (b)~$M_\odot c^2 = \num{1.9884e30}\times
(\num{2.998e8})^2 = \qty{1.787e47}{J} = \qty{1.115e57}{GeV}$, et
$\qty{1}{pc^3} = (\qty{3.0857e18}{cm})^3 = \qty{2.938e55}{cm^3}$~: le
rapport vaut \qty{37.96}{GeV.cm^{-3}}. Donc $0.30/37.96 = \num{0.0079}$ masse
solaire par parsec cube. (c)~La densité déduite de la courbe plate compte toute la masse
qui retient les orbites, étoiles et gaz compris, comme si elle était répartie en
sphère~; la matière noire n'en est qu'une partie, environ la moitié ici.
\end{solution}

\begin{solution}{exo:b3:milky-way-galaxies:7}
(a)~$W_I(\mathrm{LMC}) = -3.314\times1.4771 + 15.888 = \qty{10.993}{mag}$,
donc $M_W = 10.993 - 18.477 = \qty{-7.484}{mag}$. (b)~$\mu = 25.00 + 7.484 =
32.484$ et $d = 10^{32.484/5 + 1}~\mathrm{pc} = \qty{31.4}{Mpc}$.
(c)~$\delta d/d = (\ln 10/5)\,\delta\mu = 0.4605\times0.077 =
\qty{3.5}{\percent}$ pour une étoile, $3.5/\sqrt{50} = \qty{0.5}{\percent}$
pour cinquante, si leurs erreurs sont indépendantes (l'erreur propre à la
distance du Nuage s'ajoute à toutes de la même façon).
\end{solution}

\begin{solution}{exo:b3:milky-way-galaxies:8}
$\mu = I - A_I - M_I = 26.00 - 0.05 + 4.05 = 30.00$, donc $d =
10^{7}~\mathrm{pc} = \qty{10.0}{Mpc}$. Une erreur de \qty{0.04}{mag} correspond à
$0.4605\times0.04 = \qty{1.8}{\percent}$ sur la distance.
\end{solution}

\begin{solution}{exo:b3:milky-way-galaxies:9}
(a)~$d = cz/H_0$ donne $5\log_{10}d = 5\log_{10}(cz) - 5\log_{10}H_0$,
donc $m_B - M_B - 25 = 5\log_{10}(cz) - 5\log_{10}H_0$~; avec
$\log_{10}(cz) = a_B + 0.2\,m_B$, $5\log_{10}H_0 = 5a_B + M_B + 25$, ce qui
est le résultat. (b)~$\log_{10}H_0 = -3.8506 + 0.714158 + 5 = 1.8636$, donc
$H_0 = \qty{73.04}{km.s^{-1}.Mpc^{-1}}$. (c)~$M_B = 5(\log_{10}67.4 -
5.714158) = \qty{-19.427}{mag}$~: les supernovae sont plus brillantes de \qty{0.174}{mag},
d'où toutes les distances plus grandes d'un facteur $10^{0.2\times0.174} = 1.084$, soit le
rapport $73.04/67.4$. La tension est un désaccord de huit pour cent sur
l'\omterm{def:b3:milky-way-galaxies:ladder}{échelle des distances}.
\end{solution}

\begin{solution}{exo:b3:milky-way-galaxies:10}
(a)~Une étoile au bord, sur une orbite circulaire, vérifie $V^2/R \approx GM/R^2$
(approximation sphérique). (b)~$R = \sqrt{M/(\pi\Sigma)}$, donc $V^2 =
GM/R = G\sqrt{\pi\Sigma}\,M^{1/2}$ et $M = V^4/(\pi\Sigma G^2) \propto
V^4$~; avec un rapport masse-luminosité commun, $L \propto V^4$ aussi.
(c)~$(100/229)^4 = \num{0.036}$~: la galaxie la plus rapide est 27 fois plus
massive. (d)~À luminosité égale, le flux décroît en $1/d^2$, donc la galaxie la
moins brillante est deux fois plus éloignée.
\end{solution}

\begin{solution}{exo:b3:milky-way-galaxies:11}
(a)~$\qty{1}{km.s^{-1}} = \qty{1.0227}{kpc.Ga^{-1}}$, donc
\qty{229}{km.s^{-1}} vaut \qty{234.2}{kpc.Ga^{-1}}. (b)~$2\pi R/V$~:
\qty{107}{Ma} à \qty{4}{kpc}, \qty{322}{Ma} à \qty{12}{kpc}.
(c)~$\Delta\phi = Vt\,(1/4 - 1/12) = 234.2\times0.1667 = \qty{39.0}{rad}$,
soit 6.2 tours~: le bras serait une spirale de six tours serrés entre 4 et
\qty{12}{kpc}, bien plus serrée que les bras observés, et de plus en plus serrée à chaque
rotation supplémentaire. (d)~Dans une onde de densité, le motif tourne à une vitesse angulaire
unique, comme un corps solide~; les étoiles et le gaz la traversent, si bien que la forme
des bras ne s'enroule pas.
\end{solution}

\begin{solution}{exo:b3:milky-way-galaxies:12}
(a)~$a = (4.2\times10^6\times16^2)^{1/3} = \qty{1024}{au}$~; à
\qty{8178}{pc}, cela sous-tend $1024/8178 = \ang{;;0.125}$ (une au à un parsec
sous-tend une seconde d'arc). (b)~$R_S = 2GM/c^2 = 2\times\num{6.674e-11}
\times\num{8.35e36}/(\num{2.998e8})^2 = \qty{1.24e10}{m} = \qty{0.083}{au}$~;
le diamètre $2R_S$ sous-tend $0.166/8178 = \num{2.0e-5}$ seconde d'arc, environ
\qty{20}{\micro as}. (c)~$R_S = \qty{4.02e-7}{pc}$, et $M/(\tfrac43\pi
R_S^3) = \num{1.5e25}$ masses solaires par parsec cube, quarante fois la
borne inférieure fixée par le mouvement et la taille de Sgr~A*~: la borne se
situe à un facteur quarante près de la densité du trou noir lui-même.
\end{solution}

\begin{solution}{pb:b3:milky-way-galaxies:1}

\textbf{Partie I.}

\textbf{1.} $v_r = A\,d\sin2\ell$~: $+A d$ à \ang{45}, 0 à \ang{90}, $-A d$ à
\ang{135}, 0 à \ang{180}.

\textbf{2.} $A = (15.4 + 15.2)/(2\times\qty{1.00}{kpc}) =
\qty{15.3}{km.s^{-1}.kpc^{-1}}$.

\textbf{3.} À \ang{90}, $v_t = d(B - A) = 4.74\times(-5.91)\times1.00 =
\qty{-28.02}{km.s^{-1}}$, donc $B = 15.3 - 28.02 =
\qty{-12.72}{km.s^{-1}.kpc^{-1}}$.

\textbf{4.} À \ang{0}, $v_t = d(A + B) = 4.74\times0.55 =
\qty{2.61}{km.s^{-1}}$, donc $B = 2.61 - 15.3 = \num{-12.69}$~: les deux
concordent à \num{0.03} près.

\textbf{5.} $\omega_0 = A - B = \qty{28.0}{km.s^{-1}.kpc^{-1}}$ et
$V_0 = 28.0\times8.18 = \qty{229}{km.s^{-1}}$.

\textbf{6.} $V'_0 = -(A + B) = \qty{-2.6}{km.s^{-1}.kpc^{-1}}$~: la vitesse décroît doucement.

\textbf{7.} $P = 2\pi/\omega_0$~: $2\pi\times\qty{8.18}{kpc}/\qty{229.2}{km.s^{-1}}
= \qty{219}{Ma}$~; $4540/219 = 20.7$ tours.

\textbf{Partie II.}

\textbf{8.} Le point tangent exige une ligne de visée qui passe plus près du
centre que ne l'est le Soleil, $\sin\ell > 0$ avec $\ell < \ang{90}$ (ou le
quadrant symétrique, $\ang{270} < \ell < \ang{360}$, où la plus grande
vitesse est la plus négative). Au-delà du cercle solaire, la distance
au centre croît tout le long de la ligne de visée~: aucun point ne
se distingue, et la vitesse seule ne dit pas où se trouve le gaz.

\textbf{9.} $R = 8.18\sin\ell$~: \qtylist{5.26;6.27;7.08;7.69}{kpc}.

\textbf{10.} $V = v_\mathrm{max} + V_0\sin\ell$~: $82.9 + 147.2 = 230.1$,
$55.6 + 175.4 = 231.0$, $32.4 + 198.3 = 230.7$, $14.9 + 215.2 =
\qty{230.1}{km.s^{-1}}$.

\textbf{11.} $\Delta\nu = \nu v/c = \num{1.420406e9}\times82.9/299792 =
\qty{393}{kHz}$, vers les basses fréquences~: le gaz s'éloigne.

\textbf{12.} À environ \qty{3}{km.s^{-1}} près des géantes rouges~: entre 5 et
\qty{8}{kpc}, la courbe est plate.

\textbf{13.} Une masse ponctuelle donne $V = V_0\sqrt{R_0/R} = 229\times\sqrt{8.18/5.26} =
\qty{286}{km.s^{-1}}$ au premier point tangent, bien plus vite que la
\qty{230}{km.s^{-1}} obtenue~: la masse n'est pas concentrée au centre,
une grande partie se situe entre 5 et \qty{8}{kpc} et au-delà.

\textbf{Partie III.}

\textbf{14.} $\num{1.067e10} + 2\times\num{3.944e10} = \num{8.96e10}$ masses solaires,
presque le double de celle de la revue \num[separate-uncertainty=true]{5 \pm 1 e10}~; les étoiles du modèle
sont peut-être trop lourdes, ce qui ferait de la matière noire trouvée plus bas
une sous-estimation.

\textbf{15.} $V_b^2 = GMR^2/(R^2 + b^2)^{3/2} = \num{4.301e-6}\times\num{1.067e10}
\times400/(400.09)^{3/2}$, $V_b = \qty{47.9}{km.s^{-1}}$.

\textbf{16.} Disque mince~: $(a + b)^2 = 30.80$, $V^2 = \num{4.301e-6}\times\num{3.944e10}
\times400/(430.8)^{3/2}$, $V = \qty{87.1}{km.s^{-1}}$. \omterm{def:b3:milky-way-galaxies:structure}{Disque épais}~:
$(a + b)^2 = 11.56$, $V = \qty{90.1}{km.s^{-1}}$.

\textbf{17.} $V_\mathrm{vis} = \sqrt{47.9^2 + 87.1^2 + 90.1^2} =
\qty{134.2}{km.s^{-1}}$.

\textbf{18.} À \qty{8.18}{kpc}~: 74.8, 108.4 et \qty{127.8}{km.s^{-1}},
ensemble \qty{183.5}{km.s^{-1}}, contre $V_0 = \qty{229}{km.s^{-1}}$~:
la masse visible est déjà insuffisante au niveau du Soleil.

\textbf{19.} $\sqrt{G\times\num{8.96e10}/20} = \qty{138.8}{km.s^{-1}}$. La
valeur du modèle, \qty{134.2}{km.s^{-1}}, est proche, car à \qty{20}{kpc},
quatre à six longueurs d'échelle du disque, presque toute la masse visible se
trouve à l'intérieur de l'orbite~: la courbe visible est devenue képlérienne.

\textbf{Partie IV.}

\textbf{20.} $M = 20\times200^2/\num{4.301e-6} = \num{1.860e11}$ masses solaires.

\textbf{21.} $M_\mathrm{vis} = 20\times134.2^2/\num{4.301e-6} = \num{8.37e10}$
masses solaires.

\textbf{22.} $M_\mathrm{dark} = \num{1.860e11} - \num{8.37e10} = \num{1.02e11}$~;
puisque les deux masses valent $RV^2/G$ au même $R$, la fraction vaut
$1 - (V_\mathrm{vis}/V)^2 = 1 - (134.2/200)^2 = \num{0.550}$.

\textbf{23.} $x = 20/14.8 = 1.351$, $\ln(2.351) - 1.351/2.351 = 0.8549 - 0.5748
= 0.2801$, $M = 4\pi\times\num{1.06e7}\times14.8^3\times0.2801 =
\num{1.21e11}$ masses solaires~; la fraction de matière noire vaut
$1.21/(1.21 + 0.837) = \num{0.59}$. Le halo est ajusté à l'ensemble
de la courbe, qu'il dépasse à \qty{20}{kpc}~: les mesures y chutent
d'environ \qty{10}{km.s^{-1}} sous le modèle lissé (\qty{209}{km.s^{-1}}).

\textbf{24.} $x = 8.18/14.8 = 0.553$, $\rho = \num{1.06e7}/(0.553\times1.553^2)
= \num{7.95e6}$ masses solaires par kiloparsec cube, \num{0.0080} par
parsec cube, soit \qty{0.30}{GeV.cm^{-3}}, la valeur que donnent les auteurs de
l'ajustement.

\textbf{25.} $1 - 8.37/18.60 = \num{0.550}$~: la fraction de matière noire de la
masse contenue dans \qty{20}{kpc} est
\textbf{\qty[text-series-to-math=true]{55}{\percent}}.
\end{solution}
